\chapter{Training Experiments}


We performed the training of our policies on the NVIDIA GeForce GPUs available in our lab, and we then let the training run while randomizing the starting states in position and velocity. The initial policies were trained with a relatively myopic discount factor ($\gamma = 0.991$) and a capture radius of $0.75~\mathrm{m}$. The hyperparameters used throughout these experiments are listed in Appendix~\ref{appendix:hyperparameters}. We tested a number of different hyperparameter settings before converging on these, so their choice for the trainings in this work is empirical.


We chose to focus on three different test cases for this thesis, where we varied only the actuation limit $u_{\max}$, the maximum velocity $v_{\max}$, and the maximum initial velocity in any given episode $v_{\max}^{\mathrm{IC}}$, while keeping the remaining training hyperparameters equivalent between the three runs. We ran these with the Kalman filter off for a fully observable run, and with the Kalman filter on for a partially observable run via behavior cloning. The KF-on results have a measured-control noise standard deviation of $0.025~\mathrm{m/s^2}$ on the prior adversary action.

The three test cases that we'll focus on in the remainder of this work are:


\begin{itemize}
    \item Training case 1: $u_{\max} = 0.1~\mathrm{m/s^2}$, $v_{\max} = 1.0~\mathrm{m/s}$, $v_{\max}^{\mathrm{IC}} = 0.05~\mathrm{m/s}$
    \item Training case 2: $u_{\max} = 0.5~\mathrm{m/s^2}$, $v_{\max} = 1.5~\mathrm{m/s}$, $v_{\max}^{\mathrm{IC}} = 0.25~\mathrm{m/s}$
    \item Training case 3: $u_{\max} = 2.0~\mathrm{m/s^2}$, $v_{\max} = 1.0~\mathrm{m/s}$, $v_{\max}^{\mathrm{IC}} = 1.0~\mathrm{m/s}$
\end{itemize}

These values are meant to test the PARL+BC-produced policies and their corresponding Monte Carlo evaluations at different actuation limits and max velocity saturation limits. 

While these runs produced reasonable outcomes, certain combinations of $u_{\max}$, $v_{\max}$, and $v_{\max}^{\mathrm{IC}}$ can produce bad trainings especially at lower actuation settings. Refer to appendix \ref{appendix:failed-training-runs} if interested in seeing the behaviors that would encompass a failed run that can be seen during the training process. Furthermore, appendix \ref{appendix:legacy-ekf-bc-noise-study} provides a brief overview of the effects of increasing the prior noise on the EKF filtering across three different settings.


\section{Training Curves}
\label{sec:Net_Training_Curves}
\label{sec:training-curves}

\section{Training case 1: $u_{\max} = 0.1~\mathrm{m/s^2}$, $v_{\max} = 1.0~\mathrm{m/s}$, $v_{\max}^{\mathrm{IC}} = 0.05~\mathrm{m/s}$}
\label{sec:training-case-1}

\Needspace{0.45\textheight}
\begin{figure}[H]
\centering
\begin{minipage}[t]{0.48\textwidth}
\centering
\textbf{KF off}\\[0.25em]
\includegraphics[width=\linewidth]{\detokenize{Figures/Merged_Thesis/Training_Policy_0.1u_1vmax_0.05_icVmax/combined_training_curve_kf_off.png}}
\end{minipage}
\hfill
\begin{minipage}[t]{0.48\textwidth}
\centering
\textbf{KF on}\\[0.25em]
\includegraphics[width=\linewidth]{\detokenize{Figures/Merged_Thesis/Training_Policy_0.1u_1vmax_0.05_icVmax/combined_training_curve_kf_on.png}}
\end{minipage}
\caption{Combined training return curves for $u_{\max} = 0.1~\mathrm{m/s^2}$, $v_{\max} = 1.0~\mathrm{m/s}$, $v_{\max}^{\mathrm{IC}} = 0.05~\mathrm{m/s}$, across 5 interchanging training phases.}
\label{fig:mc-training-policy-0-1u-1vmax-0-05-icvmax-training-curves-kf}
\end{figure}


\section{Training case 2: $u_{\max} = 0.5~\mathrm{m/s^2}$, $v_{\max} = 1.5~\mathrm{m/s}$, $v_{\max}^{\mathrm{IC}} = 0.25~\mathrm{m/s}$}
\label{sec:training-case-2}

\Needspace{0.45\textheight}
\begin{figure}[H]
\centering
\begin{minipage}[t]{0.48\textwidth}
\centering
\textbf{KF off}\\[0.25em]
\includegraphics[width=\linewidth]{\detokenize{Figures/Merged_Thesis/Training_Policy_0.5u_1.5vmax_0.25_icVmax/combined_training_curve_kf_off.png}}
\end{minipage}
\hfill
\begin{minipage}[t]{0.48\textwidth}
\centering
\textbf{KF on}\\[0.25em]
\includegraphics[width=\linewidth]{\detokenize{Figures/Merged_Thesis/Training_Policy_0.5u_1.5vmax_0.25_icVmax/combined_training_curve_kf_on.png}}
\end{minipage}
\caption{Combined training return curves for $u_{\max} = 0.5~\mathrm{m/s^2}$, $v_{\max} = 1.5~\mathrm{m/s}$, $v_{\max}^{\mathrm{IC}} = 0.25~\mathrm{m/s}$, across 5 interchanging training phases.}
\label{fig:mc-training-policy-0-5u-1-5vmax-1-0-icvmax-training-curves-kf}
\end{figure}


\section{Training Case 3: $u_{\max} = 2.0~\mathrm{m/s^2}$, $v_{\max} = 1.0~\mathrm{m/s}$, $v_{\max}^{\mathrm{IC}} = 1.0~\mathrm{m/s}$}
\label{sec:training-case-3}

\Needspace{0.45\textheight}
\begin{figure}[H]
\centering
\begin{minipage}[t]{0.48\textwidth}
\centering
\textbf{KF off}\\[0.25em]
\includegraphics[width=\linewidth]{\detokenize{Figures/Merged_Thesis/Training_Policy_2.0u_1vmax_1.0_icVmax/combined_training_curve_kf_off.png}}
\end{minipage}
\hfill
\begin{minipage}[t]{0.48\textwidth}
\centering
\textbf{KF on}\\[0.25em]
\includegraphics[width=\linewidth]{\detokenize{Figures/Merged_Thesis/Training_Policy_2.0u_1vmax_1.0_icVmax/combined_training_curve_kf_on.png}}
\end{minipage}
\caption{Combined training return curves for $u_{\max} = 2.0~\mathrm{m/s^2}$, $v_{\max} = 1.0~\mathrm{m/s}$, $v_{\max}^{\mathrm{IC}} = 1.0~\mathrm{m/s}$, across 5 interchanging training phases.}
\label{fig:mc-training-policy-2-0u-1vmax-1-0-icvmax-training-curves-kf}
\end{figure}

\subsection{Training Analysis}
\label{subsec:training-analysis}

Overall, just from analyzing each curve described in Figures~\ref{fig:mc-training-policy-0-1u-1vmax-0-05-icvmax-training-curves-kf}, \ref{fig:mc-training-policy-0-5u-1-5vmax-1-0-icvmax-training-curves-kf}, and \ref{fig:mc-training-policy-2-0u-1vmax-1-0-icvmax-training-curves-kf}, we see that the PARL+BC framework manages to produce strong, convergent policies across the multiple phases of training. Furthermore, the partially observable results follow loosely the same pattern as the fully observable ones.

By studying the mean reward at which the policies converge, we see that the defender manages to get slightly stronger policies than that of the attacker as the phases progress and, consequently, the attacker policies become slightly weaker. This is a good sign, as it signals that the defender is progressively stronger and more generalizable thanks to the opponent randomization strategy described earlier, allowing it to be resilient throughout the progression of the training.

 
